\section{Рекомендации по обработке и представлению данных}

Обработка данных, как правило, осуществляется с помощью самостоятельно разработанной студентом компьютерной программы. Детальные описания методов обработки данных приведены в описаниях опытов или в пособии «Обработка данных».
Процесс обработки данных состоит из следующих этапов:
\begin{enumerate}
    \item Подготовка расчетной программы, бланков сбора данных (бумажных или электронных).
    \item Сбор данных, занесение их в бланки.
    \item Перенос данных из бумажных бланков в электронные (если требуется).
    \item Обработка данных с помощью программы.
    \item Визуализация обработанных данных в виде таблиц и графиков.
\end{enumerate}

В рамках математических моделей, описывающих исследуемые в курсе процессы и приемов, используемых для обработки описывающих эти процессы данных, составление такой программы выполняется на любом современном языке программирования высокого уровня общего назначения без использования специфических библиотек. При составлении программы следует следовать следующим рекомендациям:
\begin{enumerate}
    \item Рекомендуется использовать специализированную для выбранного языка среду разработки и следовать принятому для выбранного языка стилю кодирования.
    \item Все введенные программистом идентификаторы именуются на английском языке. Физические величины обозначаются принятыми для них буквами (если величина обозначается греческой буквой, пишется название буквы, например $\tau$ -> tau).
    \item Исходные данные считываются из внешних файлов (электронных бланков).
    \item Обработанные данные сохраняются в виде таблицы CSV (comma-separated values, «значения, разделенные запятой») или в виде фрагмента исходного кода TeX, кодирующего соответствующие таблицы и графики.
\end{enumerate}

\section{Порядок проведения работ и сдачи отчетов}
\subsection{Проведение лабораторных работ}
\begin{enumerate}
    \item Перед началом занятий надо ознакомиться и инструкциями по технике безопасности, прослушать лекцию о правилах поведения в лаборатории и расписаться в журнале.
    \item Занятия проходят по две пары подряд раз в две недели.
    \item Работы выполняются в группах по 2 (в исключительных случаях 3) человека.
    \item Тема очередной работы назначается на предыдущем занятии (или перед первым занятием).
    \item Если Вы пропустили работу или не были до нее допущены, работа выполняется в конце семестра по согласованию.
    \item Допуск до лабораторной работы состоит в собеседовании по теоретической части работы, постановочной части эксперимента и постановочной части. обработки данных. Собеседование, за исключением первого занятия, не может состояться без предъявления предварительной версии отчета по предыдущей работе (см. далее порядок сдачи отчетов).
    \item Во время собеседования можно использовать личный рукописный конспект, составленный во время подготовки к работе. Пользоваться литературой и электронными устройствами нельзя. Рекомендуется вести этот конспект в одной тетради на все лабораторные работы.
    \item Работа начинается после получения допуска и завершается после подписания протокола наблюдения преподавателем. Установку можно включать и выключать после разрешения преподавателя, ассистента или инженера.
    \item Если установка не работает или работает неадекватно, надо остановить работу и позвать преподавателя, ассистента или инженера
    \item Самовольные манипуляции с электрической сетью и попытки <<чинить>> оборудование почти наверное приведут к аннулированию протокола наблюдений: работу придется начать заново.
\end{enumerate}

\subsection{Порядок сдачи отчетов}
\begin{enumerate}
    \item Для каждой лабораторной работы создается репозиторий на \url{https://github.com} \cite{repo}, в котором должен быть исходный код отчета и всех программ для обработки данных.
    \item Репозитории должны принадлежать к одному и тому же аккаунту.
    \begin{itemize}
        \item Идеальное имя аккаунта - \verb"фамилия.и.о", например \verb"ivanov.i.i". Если имя аккаунта занято, добавьте через дефис номер: \verb"ivanov.i.i-42". \textbf{Аккаунты с непонятными именами будут проигнорированы при проверке}
        \item Репозитории должны называться \verb"Workshop1", \verb"Workshop2", \dots
    \end{itemize}
    \item Устанавливаются следующие крайние даты выгрузки информации в репозитории:
    \begin{itemize}
        \item Скан протокола наблюдений: 1 неделя после лабораторной работы.
        \item Предварительный вариант отчета и расчетные программы: 2 недели после лабораторной работы.
        \item Окончательный отчет и полностью отлаженные расчетные программы: к дате проведения зачета.
    \end{itemize}
\end{enumerate}

В любом случае предварительный отчет должен быть готов к следующему занятию.

% \subsection{Сдача отчетов в системе hwproj.ru}
% Применимо только если преподаватель дал указание использовать систему.
% \begin{enumerate}
%     \item Зарегистрироваться на портале \url{https://hwproj.ru/} используя Ваше реальное имя и почту СПбГУ. Не забудьте привязать аккаунт GitHub к созданной учетной записи
%     \item На странице \href{https://hwproj.ru/courses}{\verb"КУРСЫ"}, \verb"ВСЕ КУРСЫ" найти курс \textsc{<<Методы измерений и электромеханические системы, лабораторная работа>>}, предназначенный для Вашей академической группы, перйти на страничку курса и нажать кнопку\\ \colorbox{blue}{\color{white}\verb"ЗАПИСАТЬСЯ"}.
%     \item После обработки заявки на страничке \url{https://hwproj.ru/} появится список дедлайнов по сдаче отчетов.
%     \item Чтобы Ваш протокол наблюдений или отчет увидел преподаватель, убедитесь, что материал загружен на GitHub, найдите этот материал в списке дедлайнов, перейдите на страничку задания, нажмите кнопку\\
%     \colorbox{blue}{\color{white}\verb"ДОБАВИТЬ РЕШЕНИЕ"}. 
%     В открывшейся форме введите адрес репозитория, в котором находится материал и нажмите кнопку \colorbox{blue}{\color{white}\verb"ОТПРАВИТЬ РЕШЕНИЕ"}.
%     \item Обратите внимание, что в системе отображается порядковый номер выполнения работы, а не номер со стола в лаборатории -- он должен быть на титульной странице отчета
%\end{enumerate}
